\documentclass{beamer}
\usefonttheme{professionalfonts}
%\usetheme{Rochester}
%\usetheme{AnnArbor}
\usetheme{Szeged}
\useoutertheme{infolines}
\usecolortheme{dolphin}
\usepackage{graphicx}
% \usepackage[english]{babel}
\usepackage{multicol}
\usepackage{color}
\usepackage{xcolor}
%\usepackage{smartdiagram}


%\definecolor{carmine}{rgb}{0.59, 0.0, 0.09}
\usepackage{amsmath}

\usepackage{graphicx}
\usepackage{tikz}
\usetikzlibrary{shapes.callouts}
\usepackage{multirow}

\usepackage[utf8]{inputenc}

\usepackage{enumerate}
\usepackage[shortlabels]{enumitem}
\usepackage{amsmath,amsthm,amssymb,amsfonts,mathrsfs,latexsym,stmaryrd}
\usepackage{tcolorbox}

\usepackage{color}
\usepackage{amssymb, amsmath, amsbsy}
\usepackage{adjustbox,lipsum}
\usepackage{placeins}

\usepackage{textpos}
\usepackage{wrapfig}
\usepackage[spanish, es-tabla]{babel}

\usepackage{caption}
\setlength{\parindent}{0pt}
\hypersetup{hyperfootnotes=false}

\usepackage{url}
\usepackage{float}
\usepackage{tikz}
\usetikzlibrary{patterns}


\title[ICS-294]{Econometría ICS-294\\ \small{Clase 9: Interpretaci\'on y forma funcional}}

\author{Francisco Alfaro Medina}
\institute[USM]
{Universidad Técnica Federico Santa María\\ % Your institution for the title page
	\medskip
	\textit{francisco.alfaro@usm.cl} % Your email address

}
\logo{

\includegraphics[width=3cm]{Figuras/logoUSM.png}
}
\date{}
\newcommand\Fontvi{\fontsize{9}{7.2}\selectfont}
% Percent literal robusto: evita el shorthand \es@sppercent de babel-spanish
% que provoca "Incompatible glue units" al usar \% en modo matemático.
\newcommand{\pc}{\text{\char37}}
\begin{document}

\begin{frame}
	\titlepage % Print the title page as the first slide


\end{frame}
%%%%%%%%%%%%%%%%%%%%%%%%%%%%%%%%%%%%%%%%%%%%%%%%%%%%%%%%%%%%%%%%%%%%%%%%%%%
\logo{}

%========================= OBJETIVOS =========================
\begin{frame}{Objetivos de la Clase}

Al finalizar esta clase, el estudiante será capaz de:

\begin{itemize}
    \item[\checkmark] Interpretar correctamente los coeficientes de los modelos nivel-nivel, nivel-log, log-nivel y log-log.
    \item[\checkmark] Derivar y justificar por qué en el modelo log-lin $100\cdot\beta_1$ aproxima el cambio porcentual, y aplicar la corrección exacta $e^{\beta_1}-1$.
    \item[\checkmark] Calcular el efecto marginal en modelos con términos cuadráticos y encontrar el punto de quiebre $x^*=-\beta_1/(2\beta_2)$.
    \item[\checkmark] Interpretar términos de interacción y evaluar efectos parciales en valores relevantes.
    \item[\checkmark] Entender cómo el reescalamiento de variables afecta (y no afecta) la inferencia.
\end{itemize}

\end{frame}

\begin{frame}[shrink=10]
\frametitle{Incorporación de no-linealidades}
    \begin{enumerate}
        \item {\color{blue}Variable dependiente e independiente en niveles:}
        \begin{itemize}
            \item[] \begin{center}$y_i = \beta_0 + \beta_1 x_{1,i} + \ldots + u_i$\end{center}
            \item → \textbf{Ejemplo:} $salario_i = \beta_0 + \beta_1 educ_i + u_i$
            \item $\beta_1 = \frac{\Delta y}{\Delta x_1}$
            \item Por cada año más de educación se obtiene, en promedio, $\beta_1$ dólares más de salario.
        \end{itemize}
        \item {\color{blue}Variable dependiente en log e independiente en niveles:}
        \begin{itemize}
            \item[] \begin{center} $\log(y_i) = \beta_0 + \beta_1 x_{1,i} + \ldots + u_i$\end{center}
            \item → \textbf{Ejemplo:} $\log(salario)_i = \beta_0 + \beta_1 educ_i + u_i$
            \item $\beta_1 = \frac{\Delta \log y}{\Delta x_1}\approx \frac{\Delta \pc  y}{100\Delta x_1}$
            \item Por cada año más de educación se obtiene, en promedio, un aumento de $100 \cdot \beta_1 \pc$ del salario, manteniendo todas las demás variables constantes.
        \end{itemize}
        \item {\color{blue}Variable dependiente e independiente en logs:}
        \begin{itemize}
            \item[] \begin{center}$\log(y_i) = \beta_0 + \beta_1 \log(x_{1,i}) + \ldots + u_i$\end{center}
            \item → \textbf{Ejemplo:} $\log(salario)_i = \beta_0 + \beta_1 \log(educ)_i + u_i$
            \item $\beta_1 = \frac{\Delta \log y}{\Delta \log x_1}\approx \frac{\Delta\pc  y}{ \Delta\pc  x_1}$
            \item Un aumento de 1\% de los años de educación implica, en promedio, un aumento de $\beta_1 \pc$ del salario manteniendo todas las demás variables constantes. Es decir, representa una {\bf elasticidad}.
        \end{itemize}
    \end{enumerate}


\end{frame}

\section{Variables en logaritmo}

\begin{frame}
\frametitle{Variables en logaritmo}

\begin{itemize}
    \item Diferencias logarítmicas:
    \[
    \Delta \log(X)= \log(X_1)-\log(X_0)= \log\left (\frac{ X_1}{X_0}\right )=\log\left(1+{\Delta X\over X_0} \right)
    \]
    \item Recuerde que en nuestro curso y en el libro $log(X)=ln(X)$.
    \item Sabemos que\footnote{Serie de Taylor centrada en 0, es decir, es una buena aproximación para x cercanos a 0.}:

    \[
    \log(1+x) \approx x
    \]

\end{itemize}



\end{frame}

%========================= DERIVACIÓN APROXIMACIÓN LOG =========================
\begin{frame}[allowframebreaks]{¿Por qué $100\cdot\beta_1$ es el cambio porcentual? (log-lin)}

\begin{alertblock}{Proposición: interpretación porcentual del modelo log-lin}
En el modelo $\log(y) = \beta_0 + \beta_1 x + u$, manteniendo $u$ constante, un aumento de $\Delta x = 1$ produce un cambio porcentual en $y$ de aproximadamente $100\cdot\beta_1\,\pc$, con corrección exacta $100\cdot(e^{\beta_1}-1)\,\pc$.
\end{alertblock}

\begin{exampleblock}{Demostración}
\textbf{Paso 1.} Diferenciando $\log(y) = \beta_0 + \beta_1 x$ ante un cambio $\Delta x$ (con $u$ fijo):
\[
\Delta \log(y) = \beta_1\,\Delta x .
\]

\textbf{Paso 2.} Por definición de diferencia logarítmica, si $y$ pasa de $y_0$ a $y_1$:
\[
\Delta \log(y) = \log(y_1) - \log(y_0) = \log\!\left(\frac{y_1}{y_0}\right) = \log\!\left(1 + \frac{\Delta y}{y_0}\right).
\]

\textbf{Paso 3.} Usando la aproximación de Taylor $\log(1+z)\approx z$ para $z$ pequeño (con $z=\Delta y / y_0$):
\[
\Delta \log(y) \approx \frac{\Delta y}{y_0} = \frac{\Delta\pc  y}{100}.
\]

\textbf{Paso 4.} Igualando los pasos 1 y 3:
\[
\frac{\Delta\pc  y}{100} \approx \beta_1 \,\Delta x \;\;\Longrightarrow\;\; \Delta\pc  y \approx 100\cdot\beta_1\,\Delta x .
\]
Para $\Delta x = 1$ se obtiene $\Delta\pc  y \approx 100\cdot\beta_1$. \hfill $\square$
\end{exampleblock}

\framebreak

\begin{block}{Corrección exacta}
La aproximación $\log(1+z)\approx z$ pierde precisión cuando $\beta_1$ es grande. El cambio porcentual \emph{exacto} ante $\Delta x = 1$ se obtiene sin aproximar:
\[
\Delta \log(y) = \beta_1 \;\Longrightarrow\; \frac{y_1}{y_0} = e^{\beta_1}
\;\Longrightarrow\; \frac{\Delta y}{y_0} = e^{\beta_1} - 1 .
\]
Por lo tanto:
\[
\boxed{\;\Delta\pc  y \text{ (exacto)} = 100\cdot\left(e^{\beta_1}-1\right)\;}
\]
\end{block}

\begin{tcolorbox}[colback=blue!5!white, colframe=blue!75!black, title=Ejemplo numérico]
Si $\beta_1 = 0{,}083$:\quad aproximación $= 8{,}3\pc$; \quad exacto $= 100\cdot(e^{0{,}083}-1) = 8{,}65\pc$.\\
Si $\beta_1 = 0{,}30$:\quad aproximación $= 30\pc$; \quad exacto $= 100\cdot(e^{0{,}30}-1) = 34{,}99\pc$ (la aproximación ya falla).
\end{tcolorbox}

\end{frame}


\begin{frame}
\frametitle{Variables en logaritmo: Caso 1}


    \[\log(salario) = 0.584 + 0.083 \cdot educ\]
    \begin{itemize}
        \item Cada año adicional de educación se asocia con un incremento de 8,3\% del salario.
        \item Como el cambio porcentual es el mismo por cada año adicional de educación, el cambio en el salario en dólares ($e^{\log(salario)}$) por un año más de educación aumenta a medida que la educación lo hace.
    \end{itemize}

             \begin{figure}
        \centering
        \includegraphics[width=0.7\linewidth]{sem6/images/salario.png}
    \end{figure}

\end{frame}
\begin{frame}
\frametitle{Variables en logaritmo: Caso 2 - Elasticidad}

\begin{itemize}
    \item Cuando $x_j$ e $y$ se encuentran en logaritmo: un cambio de un 1\% en $x_j$ se asocia a un cambio de $\beta_j$\% en $y$.
    \item Por ejemplo, se estima la elasticidad de la demanda de vino en Chile.
    \[
    \log(\text{cantidad vino}) = 6 - 0,2 \log(\text{precio vino})
    \]
    \item Un aumento de un 1\% en el precio del vino corresponde a una caída del 0,2\% en la cantidad consumida de vino.
    \item Se interpreta como una elasticidad. Por ejemplo, la elasticidad de la demanda de vino con respecto al precio es 0,2.
\end{itemize}

\end{frame}

%========================= TABLA RESUMEN INTERPRETACIONES (anticipada) =========================
\begin{frame}{Las cuatro formas funcionales: resumen de interpretación}

\begin{table}[H]
    \centering
    \renewcommand{\arraystretch}{1.4}
    \begin{tabular}{l|c|c|l}
       \textbf{Modelo}  & \textbf{Dep.} & \textbf{Indep.} & \textbf{Interpretación de }$\beta_1$ \\ \hline
       Nivel-nivel  & $y$ & $x$ & $\Delta y = \beta_1\,\Delta x$ (unidades de $y$) \\
       Nivel-log  & $y$ & $\log(x)$ & $\Delta y = \dfrac{\beta_1}{100}\,(\Delta\pc  x)$ \\
       Log-nivel  & $\log(y)$ & $x$ & $\Delta\pc  y = (100\,\beta_1)\,\Delta x$ \;(semielast.) \\
       Log-log  & $\log(y)$ & $\log(x)$ & $\Delta\pc  y = \beta_1\,(\Delta\pc  x)$ \;(elasticidad) \\
    \end{tabular}
    \caption{Interpretación del coeficiente según la forma funcional.}
\end{table}

\begin{block}{Mnemotecnia}
El logaritmo convierte la variable a la que afecta en \textbf{términos porcentuales}: si $y$ está en log $\Rightarrow$ el efecto sobre $y$ es \%; si $x$ está en log $\Rightarrow$ la causa es un cambio \% en $x$.
\end{block}

\end{frame}

%========================= EJERCICIO RESUELTO 1 (logs) =========================
\begin{frame}{Ejercicio Resuelto 1: comparando formas funcionales}

\begin{block}{Ejercicio 1 (adaptado de Ejercicios/T2/Ejercicios 5 T2 PC)}
Un analista estima tres modelos de retorno a la educación con datos de la encuesta CASEN:
\begin{align*}
\text{(1)}\quad & \widehat{salario} = 10 + 2\cdot educ \\
\text{(2)}\quad & \widehat{\log(salario)} = 2{,}8 + 0{,}07\cdot educ \\
\text{(3)}\quad & \widehat{\log(salario)} = 1{,}4 + 0{,}80\cdot \log(educ)
\end{align*}
Interprete el coeficiente de $educ$ en cada caso.
\end{block}

\pause

\begin{exampleblock}{Solución}
\textbf{(1) Nivel-nivel:} un año más de educación se asocia, en promedio, a \textbf{2 unidades monetarias más} de salario (efecto absoluto constante).\\
\textbf{(2) Log-nivel (semielasticidad):} un año más de educación se asocia a un aumento de $100\cdot 0{,}07 = \mathbf{7\pc}$ del salario (exacto: $100(e^{0{,}07}-1)=7{,}25\pc$).\\
\textbf{(3) Log-log (elasticidad):} un aumento de \textbf{1\%} en los años de educación se asocia a un aumento de $\mathbf{0{,}80\pc}$ en el salario.
\end{exampleblock}

\end{frame}

\section{Formas cuadráticas}

\begin{frame}[shrink=10]
\frametitle{Formas cuadráticas}

\begin{itemize}
    \item El Modelo de Regresión Lineal se refiere a una relación lineal entre los coeficientes que acompañan a los regresores $X$. Pero las $X$ pueden tomar cualquier forma.
    \item Considere el modelo:
    \[
    y = \beta_0 + \beta_1 x + \beta_2 x^2 + u
    \]
    \item ¿Cuál es la interpretación de $\beta_1$ y $\beta_2$?
    \item \textcolor{blue}{Siempre piensen en términos de derivadas parciales con respecto a la variable explicativa.}
    \item En este caso, la predicción de $y$ es:
    \[
    \hat{y} = \hat{\beta}_0 + \hat{\beta}_1 x + \hat{\beta}_2 x^2
    \]
    Lo que implica que:
    \[
    \Delta \hat{y} \approx (\hat{\beta}_1 + 2 \hat{\beta}_2 x) \Delta x \Rightarrow \frac{\Delta \hat{y}}{\Delta x} \approx \hat{\beta}_1 + 2 \hat{\beta}_2 x
    \]
    \item Noten que ahora un cambio en $x$ genera un cambio esperado en $y$ que depende del nivel de $x$: el efecto parcial de $x$ en $y$ es variable.
    \item $\hat{\beta}_1$ es aproximadamente la pendiente de la relación entre $x$ y $y$ cuando $x = 0$.
\end{itemize}

\end{frame}

%========================= DERIVACIÓN EFECTO MARGINAL Y PUNTO DE QUIEBRE =========================
\begin{frame}{Efecto marginal y punto de quiebre $x^*$}

\begin{alertblock}{Proposición: efecto marginal en el modelo cuadrático}
En $\hat{y} = \hat{\beta}_0 + \hat{\beta}_1 x + \hat{\beta}_2 x^2$, el efecto marginal de $x$ es
\[
\frac{\partial \hat{y}}{\partial x} = \hat{\beta}_1 + 2\hat{\beta}_2 x,
\]
y existe un punto de quiebre $x^*$ donde el efecto cambia de signo, dado por $x^* = -\dfrac{\hat{\beta}_1}{2\hat{\beta}_2}$.
\end{alertblock}

\begin{exampleblock}{Demostración}
\textbf{Paso 1.} Derivamos $\hat{y}$ respecto de $x$ (tratando $x^2$ con la regla de la potencia):
\[
\frac{\partial \hat{y}}{\partial x} = \hat{\beta}_1 + 2\hat{\beta}_2 x .
\]
\textbf{Paso 2.} El efecto marginal cambia de signo donde se anula. Igualando a cero:
\[
\hat{\beta}_1 + 2\hat{\beta}_2 x^* = 0 \;\Longrightarrow\; x^* = -\frac{\hat{\beta}_1}{2\hat{\beta}_2}.
\]
\textbf{Paso 3.} Si $\hat{\beta}_2 < 0$ (y $\hat{\beta}_1>0$), $x^*$ es un \textbf{máximo}: el efecto es positivo para $x<x^*$ y negativo para $x>x^*$. Si $\hat{\beta}_2>0$, $x^*$ es un \textbf{mínimo}. \hfill $\square$
\end{exampleblock}

\end{frame}

\begin{frame}
\frametitle{Ejemplo: efecto de experiencia en salario}

\begin{itemize}
    \item Considere la siguiente ecuación:
    \[
    \hat{salario} = 3.73 + 0.298\,\text{exper} - 0.0061\,\text{exper}^2
    \]
    \item La experiencia tiene un efecto positivo en el salario pero este efecto va cayendo para niveles altos de experiencia (\(\hat{\beta}_2\) negativo).
    \item ¿Tiene sentido económico esto?
    \item Específicamente, el efecto parcial de la experiencia en el salario es:
\[
\triangle \hat{salario} / \triangle \text{exper} = 0.298 - 2 \cdot 0.0061\, \text{exper}
\]
    \item Cuando \(\beta_1\) y \(\beta_2\) tienen distinto signo, hay un valor de \(x\) a partir del cual el efecto \textit{ceteris paribus} de \(x\) en \(y\) cambia de dirección.
    \item En el ejemplo, este valor es \(x^* = \left|\frac{\beta_1}{2\beta_2}\right| = \frac{0.298}{2\cdot 0.0061} = 24.4\) años.
    \item Otros polinomios de mayor grado también pueden ser incluidos en un MRL.
\end{itemize}

\end{frame}
\begin{frame}
\frametitle{Ejemplo: experiencia y salario}


Si bien los coeficientes son lineales, la relación entre experiencia y salario deja de ser lineal.


             \begin{figure}
                 \centering
                 \includegraphics[width=0.75\linewidth]{sem6//images/sal_exp.png}
             \end{figure}



\end{frame}

%========================= EJERCICIO RESUELTO 2 (cuadrático) =========================
\begin{frame}{Ejercicio Resuelto 2: rendimientos decrecientes}

\begin{block}{Ejercicio 2 (adaptado de certamenes/Certamenes T2/Certamen 2 PC)}
Con datos de trabajadores chilenos se estima (salario en miles de CLP/hora):
\[
\widehat{salario} = 2{,}5 + 0{,}25\cdot exper - 0{,}005\cdot exper^2
\]
\textbf{(a)} ¿Cuál es el efecto marginal de un año más de experiencia? \quad \textbf{(b)} ¿En qué nivel de experiencia el salario alcanza su máximo? \quad \textbf{(c)} Interprete el signo de $\hat\beta_2$.
\end{block}

\pause

\begin{exampleblock}{Solución}
\textbf{(a)} $\dfrac{\partial \widehat{salario}}{\partial exper} = 0{,}25 - 2(0{,}005)\,exper = 0{,}25 - 0{,}01\,exper$.\\[4pt]
\textbf{(b)} Punto de quiebre: $\,exper^* = -\dfrac{\hat\beta_1}{2\hat\beta_2} = -\dfrac{0{,}25}{2(-0{,}005)} = \mathbf{25}$ \textbf{años}. Como $\hat\beta_2<0$, es un máximo.\\[4pt]
\textbf{(c)} $\hat\beta_2<0$ indica \textbf{rendimientos decrecientes}: el efecto positivo de la experiencia se atenúa y, pasados 25 años, un año adicional se asocia a un salario esperado \emph{menor}.
\end{exampleblock}

\end{frame}

\section{Términos de interacción}

\begin{frame}[shrink=10, t]
\frametitle{Términos de interacción}

\begin{itemize}
    \item En ciertas ocasiones el efecto parcial de un regresor \(x\) en la variable dependiente \(y\) puede depender del valor de otro regresor \(z\).
    \item En este caso, se incluye en la regresión un término de interacción \(x \times z\).
    \item Ejemplo: suponga el siguiente modelo para el precio de una vivienda
    \[
    \text{precio} = \beta_0 + \beta_1\text{msq} + \beta_2\text{dorm} + \beta_3\,\text{msq} \cdot \text{dorm} + u
    \]
    donde \text{msq} son los metros cuadrados de la vivienda y \text{dorm} el número de dormitorios.
    \item En este caso se cumple que
    \[
    \frac{\triangle \text{precio}}{\triangle \text{dorm}} = \beta_2 + \beta_3\,\text{msq}
    \]
    \item Si \(\beta_2 > 0\) y \(\beta_3 > 0\), la ecuación implica que un dormitorio adicional se relaciona con un incremento en el valor de la vivienda, y estos incrementos son mayores en viviendas de más metros cuadrados \text{msq}.
\end{itemize}

\end{frame}



\begin{frame}
\frametitle{Términos de interacción}

\begin{itemize}
    \item Para resumir el efecto de \text{dorm} en \text{precio} se puede evaluar \(\frac{\triangle \text{precio}}{\triangle \text{dorm}}\) en valores interesantes de \text{msq}, como su media, mediana o algún percentil extremo.
    \item Por ejemplo, sea el tamaño de la casa promedio \(\bar{\text{msq}}\), el efecto de un \text{dorm} más:
    \[
    \frac{\triangle \text{precio}}{\triangle \text{dorm}} = \beta_2 + \beta_3\bar{\text{msq}}. \quad
    \]
    \item \textbf{Cuidado con la interpretación de los efectos principales:} con la interacción, $\beta_2$ ya \emph{no} es ``el efecto de un dormitorio'', sino el efecto de un dormitorio cuando $\text{msq}=0$ (carente de sentido económico). Por eso conviene evaluar en valores relevantes.
\end{itemize}

\end{frame}

%========================= EJERCICIO RESUELTO 3 (interacción) =========================
\begin{frame}{Ejercicio Resuelto 3: interacción metros$\times$dormitorios}

\begin{block}{Ejercicio 3 (adaptado de Ejercicios/T2/Ejercicios 6)}
Para 500 viviendas se estima (precio en CLP):
\[
\widehat{\text{precio}} = 30{.}000 + 900\,\text{msq} + 5{.}000\,\text{dorm} + 60\,(\text{msq}\cdot\text{dorm})
\]
La mediana es $\widetilde{\text{msq}}=140$ y el percentil 90 es $\text{msq}_{90}=220$. ¿Cuánto sube el precio esperado al agregar un dormitorio en una vivienda mediana y en una grande?
\end{block}

\pause

\begin{exampleblock}{Solución}
El efecto marginal de un dormitorio depende del tamaño:
\[
\frac{\partial \widehat{\text{precio}}}{\partial \text{dorm}} = 5{.}000 + 60\,\text{msq}.
\]
\textbf{Vivienda mediana} ($\text{msq}=140$): $\;5{.}000 + 60(140) = \mathbf{13{.}400}$ CLP.\\
\textbf{Vivienda grande} ($\text{msq}=220$): $\;5{.}000 + 60(220) = \mathbf{18{.}200}$ CLP.\\[4pt]
\textbf{Interpretación:} un dormitorio adicional vale más en casas grandes; el signo positivo de $\hat\beta_3=60$ refleja la \textbf{complementariedad} entre espacio y número de dormitorios.
\end{exampleblock}

\end{frame}

\section{Reescalamiento de variables}

\begin{frame}
\frametitle{Reescalando variables}
¿Qué pasa si re-escalamos \(X\) o \(Y\)? Por ejemplo, para cambiar su unidad de medida?
\begin{itemize}

    \item Los coeficientes estimados, sus varianzas, estadísticos \(t\) y \(F\), cambian de manera tal que se preservan los efectos medidos y los resultados de los tests.
    \item En general, se reescalan variables por motivos estéticos, como reducir el número de ceros en los decimales de una estimación de un coeficiente.
    \item Este procedimiento es usual con variables económicas cuando los montos medidos son grandes o para transformar entre monedas de diferentes países.
\end{itemize}

\end{frame}
\begin{frame}
\frametitle{Ejemplo: reescalando \(Y\)}

\begin{itemize}
    \item Considere el siguiente modelo:
    \[
    {salario_{us\$}} = \beta_0 + \beta_1\text{educacion} + u
    \]
    donde ${salario_{us\$}}$ es medido en dólares.
    \item Quiero estimarlo en pesos a TC=1000 (i.e. \({salario_{chile}} = 1000 \cdot {salario_{us\$}}\)).
    \item Al multiplicar el modelo de los dos lados se obtiene
    \[
    1000 \cdot {salario_{us\$}} = 1000 \beta_0 + 1000 \beta_1 \text{educ} + 1000 \cdot u \quad
    \]
    \[
    {salario_{chile}} = \tilde{\beta}_0 + \tilde{\beta}_1 \text{educ} + \tilde{u} \quad
    \]
    donde \textcolor{red}{\(\tilde{\beta}_j = 1000 \beta_j\).}
    \item Interpretación. Asuma \(\beta_j = 20\)
    \begin{enumerate}
        \item Modelo original: un aumento de un año de educación implica un aumento de 20 dólares en el salario promedio.
        \item Modelo nuevo: un aumento de un año de educación implica un aumento de 20000 pesos en el salario promedio.
    \end{enumerate}
\end{itemize}
\text{}
\end{frame}


\begin{frame}{Reescalando un regresor X }
            \[ y = \beta_0 + \beta_1 x + u\]
Si re-escala la variable $x$ por $\tilde{x} = \alpha x$ se estima:
            \[ y = \beta_0 + \tilde{\beta}_1 \tilde{x} + u\]
        $\Rightarrow$
        \begin{itemize}
            \item $\hat{\tilde{\beta}}_1 = \frac{\hat{\beta}_1}{\alpha}$
            \item $\text{Var}(\hat{\tilde{\beta}}_1) = \dfrac{\text{Var}(\hat{\beta}_1)}{\alpha^2}$
            \item El estadístico $t$ \textbf{no} cambia: $\;t = \dfrac{\hat{\tilde\beta}_1}{\text{ee}(\hat{\tilde\beta}_1)} = \dfrac{\hat\beta_1/\alpha}{\text{ee}(\hat\beta_1)/\alpha}$.
            \item $\hat{\sigma}^2$, SSR, SST, $R^2$ \textbf{no} cambian (no se altera el ajuste).
        \end{itemize}
\end{frame}

%========================= EJEMPLO APLICADO CON CÓDIGO =========================
\begin{frame}[fragile]{Ejemplo aplicado: estimación de formas funcionales en R}

\begin{exampleblock}{Retornos a la educación con datos chilenos (CASEN)}
\begin{verbatim}
datos <- read.csv("casen.csv")

# log-lin (semielasticidad) + cuadrático en experiencia
mod1 <- lm(log(salario) ~ educ + exper + I(exper^2),
           data = datos)
summary(mod1)

# Punto de quiebre de la experiencia: -b1 / (2*b2)
b <- coef(mod1)
xstar <- -b["exper"] / (2 * b["I(exper^2)"])

# Interaccion metros cuadrados x dormitorios (vivienda)
mod2 <- lm(precio ~ msq + dorm + msq:dorm, data = viv)
\end{verbatim}
\end{exampleblock}

\footnotesize \textit{En Python: \texttt{smf.ols("np.log(salario) \textasciitilde\ educ + exper + I(exper**2)", data=datos).fit()}}

\end{frame}

%========================= EJERCICIOS PROPUESTOS =========================
\begin{frame}{Ejercicio Propuesto 1 (rutinario)}

\begin{block}{Ejercicio 1 (Propuesto)}
Se estima la elasticidad de las ventas respecto a la publicidad:
\[
\widehat{\log(\text{ventas})} = 1{,}2 + 0{,}45\,\log(\text{publicidad}).
\]
\textbf{(a)} Interprete el coeficiente $0{,}45$. \quad \textbf{(b)} Si la publicidad aumenta 10\%, ¿cuál es el cambio porcentual esperado en las ventas?
\end{block}

{\small \textit{Indicación: en un modelo log-log el coeficiente es directamente una elasticidad.}}

\end{frame}

\begin{frame}{Ejercicio Propuesto 2 (intermedio)}

\begin{block}{Ejercicio 2 (Propuesto)}
Con datos de un test de admisión se estima
\[
\widehat{\text{Puntaje}} = 200 + 10\,\text{IQ} - 0{,}05\,\text{IQ}^2 .
\]
\textbf{(a)} Halle el efecto marginal de IQ. \quad \textbf{(b)} ¿En qué valor de IQ el puntaje es máximo? \quad \textbf{(c)} ¿Qué ocurre con el efecto marginal para IQ $>$ ese valor?
\end{block}

{\small \textit{Indicación: derive respecto de IQ e iguale a cero; use $x^*=-\beta_1/(2\beta_2)$.}}

\end{frame}

\begin{frame}{Ejercicio Propuesto 3 (desafiante)}

\begin{block}{Ejercicio 3 (Propuesto)}
En el modelo $\log(salario) = \beta_0 + \beta_1\,educ + \beta_2\,D + \beta_3\,(D\cdot educ) + u$, donde $D=1$ si el trabajador es mujer:
\textbf{(a)} Escriba el retorno (semielasticidad) a la educación para hombres y para mujeres.
\textbf{(b)} ¿Qué hipótesis nula contrastaría para afirmar que el retorno a la educación es igual entre géneros?
\textbf{(c)} Si se reescala $educ$ a meses ($\tilde{educ}=12\cdot educ$), ¿cómo cambia $\hat\beta_1$, su $t$ y el $R^2$?
\end{block}

{\small \textit{Indicación: derive $\partial\log(salario)/\partial educ$ por grupo; recuerde el efecto del reescalamiento sobre $t$ y $R^2$.}}

\end{frame}

%========================= TABLA RESUMEN FINAL =========================
\begin{frame}{En resumen}

\begin{table}[H]
    \centering
    \renewcommand{\arraystretch}{1.3}
    \begin{tabular}{l|c|c|c}
       Modelo  & Var. dependiente & Var. independiente & Interpretación $\beta$ \\ \hline
       Nivel-nivel  & $y$ & $x$ & $\Delta y = \beta_i \Delta x$\\
       Nivel-log  & $y$ & $log(x)$ & $\Delta y = \frac{\beta_i}{100} \Delta \pc  x$ \\
       Log-nivel  & $log(y)$ & $x$ & $ \Delta\pc  y = (100 \beta_i) \Delta x$\\
       Log-log  & $log(y)$ & $log(x)$ & $ \Delta \pc  y = \beta_i \Delta \pc x$\\
    \end{tabular}
    \caption{Interpretación}
    \label{tab:my_label}
\end{table}

\end{frame}

%========================= RESUMEN FINAL =========================
\begin{frame}{Resumen}

\begin{block}{Puntos Clave de la Clase 9}
\begin{enumerate}
    \item El logaritmo transforma a interpretación \textbf{porcentual}: log en $y\Rightarrow$ efecto en \%; log en $x\Rightarrow$ causa en \% (log-log $=$ elasticidad).
    \item En log-lin, $100\cdot\beta_1$ aproxima el cambio \%; la versión exacta es $100\,(e^{\beta_1}-1)$, relevante cuando $\beta_1$ es grande.
    \item En modelos cuadráticos el efecto marginal es $\beta_1+2\beta_2 x$ y el punto de quiebre es $x^*=-\beta_1/(2\beta_2)$.
    \item Con interacciones, el efecto parcial de una variable depende del nivel de otra: evalúe en media/mediana/percentil.
    \item Reescalar variables cambia magnitudes de coeficientes y varianzas, pero \textbf{no} altera los estadísticos $t$, $F$ ni el $R^2$.
\end{enumerate}
\end{block}

\begin{alertblock}{Para la próxima clase}
Revisar: pruebas de hipótesis múltiples (test $F$) y el uso de variables artificiales (dummies).
\end{alertblock}

\end{frame}

\end{document}
